\section{Exploration Algorithm}\label{sec:exploration}

\subsection{Formalisation}\label{sec:exploration:formalization}

The exploration loop maintains a tree $\Tree=(\Nodes, E)$ in which
each node $n\in\Nodes$ represents a single research step: an idea
revision, a code change, a benchmark run, an analysis, etc. A node
carries

\begin{itemize}
  \item a discrete state $\statefn{n}\in\{\StateOpen,\StateEval,\StateDone,\StateFailed\}$,
  \item a categorical \term{label} drawn from the \texttt{NodeLabel}
        enumeration,
  \item per-axis evaluation metrics in a freeform dict
        \texttt{node.metrics}, including a distinguished scalar
        $\texttt{\_scientific\_score}\in[0,1]$,
  \item a Boolean \texttt{has\_real\_data} that flags whether the
        node carries actual measurements rather than provisional
        text, and
  \item the bookkeeping field \texttt{depth} surfaced both in the
        selection prompts and to the pruning predicate. ARI does not
        retry failed nodes, so no \texttt{retry\_count} signal is
        surfaced.
\end{itemize}

The aggregation of axis-level signals into the scientific score is
delegated to any implementation of the \texttt{Evaluator} protocol;
ARI does not pin a fixed linear combination. The protocol is the only seam between BFTS
and downstream rubric evaluation, which keeps the search engine
agnostic about how axes get reduced to a scalar.

\paragraph{Run-loop state.}
Each iteration of the outer loop transforms the tuple
\[
  \state \;=\; (\Frontier,\;\Pending,\;\expandcount,\;\Hist),
\]
where $\Frontier\subseteq\Nodes$ holds completed nodes
(\StateDone{} / \StateFailed{}) eligible for expansion, $\Pending\subseteq\Nodes$
holds nodes in state \StateOpen{} ready to execute, $\expandcount:\Nodes\to\Naturals$
counts how many times each frontier node has been expanded, and
$\Hist=(l_1,\ldots,l_t)$ is the sliding window of recently-executed
labels (length capped at $\histlen=20$, see~\cref{sec:exploration:diversity}).
The initial state is $\state_0=(\emptyset,\{n_\mathrm{root}\},\mathbf{0},())$.

\paragraph{Pruning predicate.}
The hard exploration budget is the predicate
\begin{multline}\label{eq:prune}
  \prunepred(n;\,|\Nodes|) \;=\;
    \indicator{|\Nodes| \geq B_N}
    \;\lor\;
    \indicator{\depth{n} \geq B_D}
    \\\;\lor\;
    \sterilepred(n),
\end{multline}
where $B_N = \texttt{config.max\_total\_nodes}$, $B_D =
\texttt{config.max\_depth}$, and the \emph{sterile} predicate
\begin{equation}\label{eq:sterile}
  \sterilepred(n) \;=\; \indicator{\,|\Delta_n| = 0\,}
\end{equation}
(with $\Delta_n$ the set of files added, modified, or deleted by $n$)
fires when a node finishes its agent loop without writing any new
artefact relative to its parent. The run loop computes this
file diff and records the result as a Boolean \texttt{\_sterile} flag
on the node; \texttt{should\_prune} reads that flag together with the
two caps. The caller passes the live $|\Nodes|$ each iteration so
there is exactly one source of truth for the node count.
Step / cost budgets are enforced separately by the cost tracker at
the call site, not inside the BFTS engine.

\paragraph{Retire predicate.}
On top of $\prunepred$, the run-loop applies a softer
\emph{frontier retire predicate} $\retirepred$ to drop a parent $f$
from $\Frontier$ once it is no longer the most productive lead:
\begin{equation}\label{eq:retire}
  \retirepred(f, c) \;=\;
    \underbrace{\indicator{\nodescore{c} > \nodescore{f}}}_{\text{Rule A}}
    \;\lor\;
    \underbrace{\indicator{\expandcount(f) \geq \maxexp}}_{\text{Rule B}},
\end{equation}
with $\maxexp = \texttt{config.max\_expansions\_per\_node}$ (default
$4$). Rule~A retires $f$ as soon as a child $c$ outscores it; Rule~B
caps the per-parent budget so the search spreads.
$\retirepred$ does \emph{not} delete the node — its history stays on
$\Tree$ — only its eligibility to be re-expanded is revoked.

The full step decomposition of one outer iteration (inner expansion
sub-loop, parallel ReAct batch, post-execution retire) is shown in
\cref{fig:bfts-state}; the explicit pseudocode is
\cref{alg:bfts-run-loop} in \cref{sec:exploration:run-loop}.

\begin{figure}[t]
  \centering
  \resizebox{0.85\columnwidth}{!}{% F15: BFTS state-machine view of one outer run-loop iteration, set in the
% ARI house design language (_style.tex). Single top-down spine; the two
% iteration scopes are dashed ariloop enclosures — the inner expansion
% sub-loop (steps 1-3) and the outer parallel-run loop (steps 4-7) —
% joined through the white state pools they communicate over (frontier
% F at the top, pending P between them). Pale-green boxes are the
% predicate gates (Pi, sigma, rho); the vermilion edges (the single
% accent role of this figure) are the side effects fired when a gate
% predicate evaluates true; the dotted gray margin path is the reflux
% into the next outer iteration. Auxiliary state (e, H) stays in the
% caption, exactly as before; per-edge state-bump annotations were
% folded into the caption to keep annotation density low.
%
% style.tikzstyles is loaded by the preamble (inlined) and by the
% standalone wrapper; the refined ari* styles used below come from
% _style.tex in this directory (gate T2: all widths live in the style
% files, never inline).
%
% The build cwd differs per entry point, so resolve _style.tex from each
% known prefix: report/<lang>/ (the language PDFs \input figures as
% ../shared/figures/tikz/<fig>.tex), report/ (make figure / render_tikz),
% the repo root, then this directory itself (direct cwd or TEXINPUTS).
\InputIfFileExists{../shared/figures/tikz/_style.tex}{}{%
  \InputIfFileExists{shared/figures/tikz/_style.tex}{}{%
    \InputIfFileExists{report/shared/figures/tikz/_style.tex}{}{%
      \InputIfFileExists{_style.tex}{}{%
        \errmessage{F15: cannot locate figures/tikz/_style.tex}}}}}
\begin{tikzpicture}[aricanvas]
  %% Heading: the state tuple the run loop mutates.
  \node[ariphase] (hd) {\figlabel{F15.band.state}};
  \pic at ([xshift=-2.5mm]hd.west) {ariglyph flask};

  %% Spine. White boxes are shared state pools, peach boxes the
  %% experimentation-phase steps, pale-green boxes the predicate gates.
  \node[aribox, below=0.4cm of hd]                    (frontier) {\figlabel{F15.frontier}};
  \node[ariboxwide, fill=phaseexp, below=0.9cm of frontier] (pick) {\figlabel{F15.pick_best}};
  \node[aribox gate, below=of pick]                   (prune)    {\figlabel{F15.prune}};
  \node[aribox exp, below=of prune]                   (expand)   {\figlabel{F15.expand}};
  \node[aribox, below=0.8cm of expand]                (pending)  {\figlabel{F15.pending}};
  \node[ariboxwide, fill=phaseexp, below=0.9cm of pending]  (batch) {\figlabel{F15.select_run}};
  \node[ariboxwide, fill=phaseexp, below=of batch]    (react)    {\figlabel{F15.run_parallel}};
  \node[aribox gate, below=of react]                  (sterile)  {\figlabel{F15.sterile}};
  \node[aribox, below=of sterile]                     (append)   {\figlabel{F15.append_frontier}};
  \node[aribox gate, below=of append]                 (retire)   {\figlabel{F15.retire}};

  %% Gate side effects (fire when the predicate is true).
  \node[aribox, right=of prune]   (dropf) {\figlabel{F15.prune_drop}};
  \node[aribox, right=of sterile] (clamp) {\figlabel{F15.clamp}};
  \node[aribox, right=of retire]  (dropr) {\figlabel{F15.retire_drop}};

  %% Loop enclosures: inner expansion sub-loop / outer parallel-run loop.
  \node[ariloop, fit=(pick)(prune)(expand)]                  (loopin)  {};
  \node[ariloop, fit=(batch)(react)(sterile)(append)(retire)] (loopout) {};
  \node[arinote, anchor=south west] at ([xshift=6mm]loopin.north)  (ttlin)  {\figlabel{F15.lane.inner}};
  \node[arinote, anchor=south west] at ([xshift=6mm]loopout.north) (ttlout) {\figlabel{F15.lane.outer}};
  \pic at ([xshift=-2.5mm]ttlin.west)  {ariglyph loop};
  \pic at ([xshift=-2.5mm]ttlout.west) {ariglyph loop};

  %% Step badges 1..7 (macro flow order; replaces edge labels).
  \node[aribadge] at (pick.north west)    {1};
  \node[aribadge] at (prune.north west)   {2};
  \node[aribadge] at (expand.north west)  {3};
  \node[aribadge] at (batch.north west)   {4};
  \node[aribadge] at (react.north west)   {5};
  \node[aribadge] at (sterile.north west) {6};
  \node[aribadge] at (retire.north west)  {7};

  %% Main top-down flow.
  \draw[ariarrow] (frontier) -- (pick);
  \draw[ariarrow] (pick)     -- (prune);
  \draw[ariarrow] (prune)    -- (expand);
  \draw[ariarrow] (expand)   -- (pending);
  \draw[ariarrow] (pending)  -- (batch);
  \draw[ariarrow] (batch)    -- (react);
  \draw[ariarrow] (react)    -- (sterile);
  \draw[ariarrow] (sterile)  -- (append);
  \draw[ariarrow] (append)   -- (retire);

  %% Gate-fired side effects: the one accent role in this figure.
  \draw[ariarrow accent] (prune)   -- (dropf);
  \draw[ariarrow accent] (sterile) -- (clamp);
  \draw[ariarrow accent] (retire)  -- node[arinote, above] {\figlabel{F15.yes_ab}} (dropr);

  %% Reflux: retire's "no" arm starts the next outer iteration.
  \draw[arifeedback] (retire.west) -- ([xshift=-1.6cm]retire.west)
                     |- (frontier.west);
  \node[arinote] at ([xshift=-1.6cm]pending.west) {\figlabel{F15.no_next}};
\end{tikzpicture}
}
  \caption{State-machine view of one outer BFTS iteration as a single
           top-to-bottom flowchart. Steps 1--3 are the inner expansion
           sub-loop; steps 4--7 are the outer parallel-run loop. Yellow
           diamonds are the predicate guards ($\prunepred,\sterilepred,\retirepred$);
           red dashed boxes on the right are the side-effects each
           predicate triggers (frontier drop / score clamp / retire).
           Auxiliary state not drawn here: the per-parent expansion
           counter $\expandcount(\cdot)$ is bumped at step~3 and
           consulted by $\retirepred$ at step~7, and the label history
           $\Hist$ (window $\histlen$) is appended at step~5 via
           \texttt{record\_run} and consulted at step~4 to compute the
           diversity bonus $\divbonus{\cdot}$.
           Compare with~\cref{fig:bftsflow}, which shows the same loop
           as a finer-grained control-flow spine with the LLM prompts
           drawn in.}
  \label{fig:bfts-state}
\end{figure}

\subsection{Run-loop algorithm}\label{sec:exploration:run-loop}

\Cref{alg:bfts-run-loop} writes the iteration as explicit pseudocode.
The inner \textsc{while} fills empty worker slots one expansion at a
time (so workers always start running before the next expansion is
proposed); the outer block then assembles a batch by calling
\texttt{select\_next\_node} repeatedly --- one node per call --- until
the batch reaches the worker cap, and executes that batch in parallel;
the post-execution block (a) records the executed label into $\Hist$
so the diversity bonus reflects what actually ran, and (b) applies
$\retirepred$ Rules~A and~B.

\begin{algorithm}[t]
\small
\caption{BFTS run-loop (one outer iteration).}
\label{alg:bfts-run-loop}
\begin{algorithmic}[1]
\Require state $\state=(\Frontier,\Pending,\expandcount,\Hist)$;
         config $(B_N, B_D, \maxexp)$ and worker cap
         $\texttt{max\_parallel}=\min(\texttt{max\_parallel\_nodes},4)$
\Ensure  updated state $\state'$
\State \Comment{--- inner expansion sub-loop ---}
\While{$\Frontier \neq \emptyset \land |\Pending| < \texttt{max\_parallel} \land |\Nodes| < B_N$}
  \State $f \gets \selectLLM(\Frontier)$ \Comment{\texttt{select\_best\_to\_expand}}
  \If{$\prunepred(f;\,|\Nodes|)$}
    \State $\Frontier \gets \Frontier \setminus \{f\}$;\;\;\textbf{continue}
  \EndIf
  \State $c \gets \texttt{expand}(f,\,\texttt{depth\_note},\,\texttt{budget\_note})$
  \State $\Pending \gets \Pending \cup \{c\}$;\;\;$\expandcount(f) \mathrel{+}{=} 1$
\EndWhile
\State \Comment{--- outer parallel-execution batch ---}
\State $B \gets \emptyset$
\While{$|B| < \min(\texttt{max\_parallel},\,|\Pending|)$}
  \State $n \gets \selectLLM(\Pending)$ \Comment{\texttt{select\_next\_node}: one node}
  \State $\Pending \gets \Pending \setminus \{n\}$;\;\;$B \gets B \cup \{n\}$
\EndWhile
\ForAll{$n \in B$ \textbf{in parallel}}
  \State $n' \gets \mathrm{ReAct}(n)$ \Comment{\texttt{AgentLoop.run}, may fail}
  \State $\Hist \gets (\Hist \mathbin{\Vert} \mathrm{label}(n'))[\,-\histlen{:}\,]$ \Comment{\texttt{record\_run}}
  \State $\Frontier \gets \Frontier \cup \{n'\}$
  \ForAll{$p \in \Frontier : p = \parentof{n'}$} \Comment{Rule A}
    \If{$\nodescore{n'} > \nodescore{p}$}
      \State $\Frontier \gets \Frontier \setminus \{p\}$
    \EndIf
  \EndFor
  \ForAll{$p \in \Frontier$} \Comment{Rule B}
    \If{$\expandcount(p) \geq \maxexp$}
      \State $\Frontier \gets \Frontier \setminus \{p\}$
    \EndIf
  \EndFor
\EndFor
\State \Return $\state' = (\Frontier,\Pending,\expandcount,\Hist)$
\end{algorithmic}
\end{algorithm}

\subsection{Node selection (LLM-driven)}\label{sec:exploration:node-selection}

\begin{figure}[t]
  \centering
  \resizebox{\columnwidth}{!}{% F03: signals consumed by the LLM-driven node selector.
% Replaces the earlier closed-form utility figure: ARI does not pin
% a scalar utility; instead a structured candidate description is
% handed to the LLM and an index is parsed back. See F11 for the full
% control flow; this figure focuses on the *information* that flows in.
% style.tikzstyles is loaded by shared/preamble.tex / standalone wrapper.
\begin{tikzpicture}[node distance=0.5cm and 1.6cm]

  % per-node signals (left column), all stations for paper-style typography
  \node[station]                                       (real)   {\figlabel{F3.has_real}};
  \node[station, below=of real]                        (metrics){\figlabel{F3.metrics}};
  \node[station, below=of metrics]                     (depth)  {\figlabel{F3.depth_signal}};
  \node[station, below=of depth]                       (label_st){\figlabel{F3.retry}};
  \node[station, below=of label_st]                    (summ)   {\figlabel{F3.summary}};
  \node[station, fill=ari-pink!22, below=of summ]      (div)    {\figlabel{F3.div}};

  % LLM selector (centre) — wider station emph
  \node[station emph big, right=2.2cm of metrics, yshift=-1.5cm]
                                                       (sel)    {\figlabel{F3.selector}};
  % the consumed prompt file (above the selector)
  \node[station, fill=ari-yellow!30, draw=ari-orange,
        above=0.6cm of sel]                            (prom)   {\figlabel{F3.prompt}};

  % output (right) — single index back to caller
  \node[station emph big, right=2.2cm of sel]          (idx)    {\figlabel{F3.index}};

  % deterministic fallback when LLM index is unparseable
  \node[station, dashed, draw=ari-vermilion, below=1.0cm of sel]
                                                       (fb)     {\figlabel{F3.fallback}};

  % --- numbered step badges --------------------------------------
  \node[numbered step, above left=0.05cm and -0.2cm of real.north west] {1};
  \node[numbered step, above left=0.05cm and -0.2cm of sel.north west]  {2};
  \node[numbered step, above left=0.05cm and -0.2cm of idx.north west]  {3};

  % --- edges -----------------------------------------------------
  \foreach \n in {real,metrics,depth,label_st,summ}
    \draw[arrow] (\n.east) -- (sel.west);
  \draw[arrow dashed] (div.east) to[bend right=8] (sel.west);
  \draw[arrow dashed] (prom.south) -- (sel.north);
  \draw[big arrow] (sel.east) -- (idx.west)
       node[edge label, font=\sffamily\footnotesize]{\figlabel{F3.parse}};
  \draw[big arrow, draw=ari-vermilion] (sel.south) -- (fb.north)
       node[edge label, font=\sffamily\footnotesize]{\figlabel{F3.unparseable}};
  \draw[big arrow, draw=ari-vermilion] (fb.east) to[bend right=22] (idx.south);

  % --- background grouping for input column ----------------------
  \begin{scope}[on background layer]
    \node[stage container, fit=(real)(metrics)(depth)(label_st)(summ)(div)] (col) {};
    \node[stage title, above=4pt of col.north west, anchor=south west] {\figlabel{F3.band.signals}};
  \end{scope}
\end{tikzpicture}
}
  \caption{Selection signals fed to the LLM that picks the next node
           to operate on. Inputs are concatenated into a structured
           candidate description; the soft \texttt{diversity\_bonus}
           is a tiebreaker, not a primary signal.}
  \label{fig:nodescore}
\end{figure}

ARI deliberately does \emph{not} reduce node ranking to a closed-form
scalar utility. Instead, both \texttt{select\_next\_node}
(``which pending node should the next agent step run?'') and
\texttt{select\_best\_to\_expand} (``which completed node is most
worth expanding?'') hand a list of candidate descriptions to the LLM
and parse back a single 0-based index --- each call selects exactly
one node. The candidate description for each node $n$ is a one-liner
of the form
\begin{quote}\small\texttt{[i] id=..., depth=..., label=..., has\_real\_data=..., metrics=\{...\}, summary=...}\end{quote}
The run-selection variant (\texttt{select\_next\_node}) additionally
appends \texttt{diversity\_bonus=+0.05} to the line for any node that
currently earns the bonus; \texttt{select\_best\_to\_expand} omits it.
The prompts that consume the list are reproduced verbatim as
\cref{prompt:bfts_select} (selection) and \cref{prompt:bfts_expand_select}
(expansion target). The selection criteria the prompt enumerates are:
nodes with \texttt{has\_real\_data=True} and strong metrics are
high-value; metrics are weighed holistically (multi-objective);
deeper nodes with excellent results are worth continuing; the
\texttt{diversity\_bonus} is treated as a soft tiebreaker only.
The \texttt{label} field matches the information shown to
\texttt{select\_best\_to\_expand}, and there is no spurious
``low retry'' criterion.

\subsubsection*{Diversity bonus}\label{sec:exploration:diversity}

The diversity bonus \texttt{BFTS.diversity\_bonus} tracks recent
labels in \texttt{\_recent\_label\_history} (a sliding window populated
by \texttt{record\_run}). For a node whose label is \emph{under-}
represented relative to the most common label in that window, the
function returns
\begin{equation}\label{eq:diversity}
  \mathrm{div}(n) =
    \begin{cases}
      +0.05 & \mathrm{cnt}(n) \le \tfrac{1}{2}\,\mathrm{cnt}_{\max},\\
      0     & \text{otherwise,}
    \end{cases}
\end{equation}
where $\mathrm{cnt}(n)$ is the number of times $\mathrm{label}(n)$
occurs in the sliding window and
$\mathrm{cnt}_{\max}=\max_\ell\mathrm{cnt}(\ell)$.
The constant $0.05$ is small by design: scientific score still
dominates ranking, and the bonus only breaks near-ties.

\subsubsection*{LLM fallback}\label{sec:exploration:fallback}

If the LLM returns an unparseable index, \texttt{select\_next\_node}
falls back to a deterministic ranking via the
\texttt{\_select\_fallback} helper (which ranks candidates by the
\texttt{\_fallback\_score} expression). The default scoring expression
is
\begin{equation}\label{eq:fallback}
  \fallbackscore{n} \;=\; \nodescore{n} + \mathrm{div}(n),
\end{equation}
where $\nodescore{n}\equiv\texttt{\_scientific\_score}(n)$ is the
node's scalar reward (\cref{sec:notation}).
\texttt{\_select\_fallback} prefers nodes with
\texttt{has\_real\_data=True} when any are available, then returns the
candidate with the highest $\fallbackscore{\cdot}$. This guarantees
the loop always makes progress even when the LLM call degrades.

\subsubsection*{Configurable evaluation layers}\label{sec:exploration:config-layers}

The four evaluation surfaces in \cref{sec:exploration:formalization,sec:exploration:node-selection,sec:exploration:fallback}
are independently configurable; the defaults reproduce the equations
above, so an unmodified configuration is a no-op.

\begin{description}
  \item[\textbf{Composite formula.}]
        Selected by \code{evaluator.composite}. The reduction
        $\{a_k\}\!\to\!\texttt{\_scientific\_score}$ is one of:
        weighted harmonic mean (default), arithmetic mean, bottleneck
        \code{weighted\_min}, or weighted geometric mean. The
        harmonic and geometric forms apply a floor $\epsilon=0.01$
        to a zero-valued axis so the denominator stays finite.
  \item[\textbf{Frontier-score strategy.}]
        Selected by \code{bfts.frontier\_score}.
        \cref{eq:fallback} is the
        \code{scientific\_plus\_diversity} case. The alternatives
        are \code{scientific\_only} (drops $\mathrm{div}$);
        \code{depth\_penalized}, adding
        $-\lambda\,\depth{n}$ on top of $\mathrm{div}(n)$
        with $\lambda=$\code{depth\_penalty\_lambda}; and
        \code{ucb\_like}, adding
        $c_{\mathrm{ucb}}\sqrt{\log N / (\expandcount(n)+1)}$ on top of
        $\mathrm{div}(n)$, where $c_{\mathrm{ucb}}=$\code{ucb\_c} and
        $N = \sum_{n'}\expandcount(n')+|\Frontier|$.
  \item[\textbf{Axis set.}]
        Selected by \code{evaluator.axis\_mode}.
        \texttt{dynamic} (default) builds the axis set from the
        generic 5-axis floor plus the active rubric and
        \texttt{idea.json} plan keywords. \texttt{legacy} pins to the
        5 canonical axes. \texttt{custom} feeds an explicit
        $\{(\mathrm{name},\mathrm{description},w)\}$ list verbatim
        through to the judge LLM.
  \item[\textbf{Selection prompts.}]
        Set via \code{bfts.select\_prompt} and
        \code{bfts.expand\_select\_prompt}: a pair of
        \code{FilesystemPromptLoader} keys lets the user
        swap in alternative templates for the two
        $\selectLLM$ calls in \cref{alg:bfts-run-loop} without
        editing source code. The selection template must accept the
        placeholders \code{experiment\_goal}, \code{memory\_context},
        and \code{candidates}; the expansion-target template accepts
        \code{experiment\_goal} and \code{candidates}.
\end{description}

\subsection{Expansion (one child per call)}\label{sec:exploration:expansion}

The \texttt{expand} function generates \emph{at most one} new child
per invocation (no pre-batching). Callers wanting more children for
the same parent call \texttt{expand} repeatedly, threading the
previously created children through the \texttt{existing\_children}
argument so the LLM can avoid proposing duplicate directions.

The label of each new child is decided by the LLM rather than by a
hardcoded template; the expansion prompt, reproduced verbatim as
\cref{prompt:bfts_expand}, is fed:

\begin{itemize}
  \item the parent's status (\texttt{succeeded} /
        \texttt{failed/no-real-data}) and
        \texttt{\_scientific\_score};
  \item the parent's structured \texttt{node\_report}
        (\texttt{delta\_vs\_parent} / \texttt{concerns} /
        \texttt{hints}, produced by the node-report builder) when
        present, so the planner can target specific weaknesses;
  \item sibling scores at the same depth, and ancestor scores from
        root to parent;
  \item tree-wide diversity metrics (unique labels seen so far,
        depth distribution); and
  \item the labels of children already spawned at this parent, so
        duplicates are not proposed.
\end{itemize}

A child is created in state \StateOpen{}, handed to the agent loop
for execution (\cref{sec:exploration:react}), and transitions to
\StateDone{} when all axes are populated or to \StateFailed{} if the
agent loop raises.

\subsection{Pruning and termination}\label{sec:exploration:termination}

\texttt{should\_prune} implements the pruning predicate $\prunepred$
(\cref{eq:prune}): total-node cap, depth cap, and the sterile flag.
The depth check activates the \texttt{max\_depth}
config field. The frontier retire predicate $\retirepred$
(\cref{eq:retire}) is applied after each node completes: Rule~A
retires a parent $f$ once a child's reward $\nodescore{\cdot}$
surpasses it, and Rule~B retires $f$ once it has been expanded
$\maxexp$ times. $\retirepred$ does not delete the node — its
history stays on $\Tree$ — but the frontier pool shrinks so the
search spreads rather than mining the same parent indefinitely.

The loop terminates when any of:

\begin{itemize}
  \item the global \texttt{max\_total\_nodes} cap is hit;
  \item \emph{every} frontier node fails $\Pi$ (depth cap, sterile,
        or budget exhausted) and the pending pool is empty;
  \item the per-call step / cost budgets enforced by the cost
        tracker run out;
  \item the \term{stagnation detector} \texttt{detect\_stagnation}
        observes that the running max of \texttt{\_scientific\_score}
        has not improved over a sliding window;
  \item a \texttt{LineageDecision} from the lineage-decision module
        explicitly halts the branch (see \cref{prompt:lineage_decision}).
\end{itemize}

Which of these fires most often is an empirical question that
\cref{chapter:eval} is meant to answer once a frozen checkpoint
is in place; the present report does not claim a distribution.

\subsection{Lineage and reproducibility}\label{sec:exploration:lineage}

A lineage is the chain of ancestors of the best node, written into
the checkpoint as the \texttt{lineage} key in \texttt{tree.json}.
The \texttt{LineageState} dataclass captures the running statistics
that BFTS uses to decide whether to continue, and the cross-checkpoint
walker \texttt{walk\_ancestor\_ckpts} provides parent pointers across
runs. Idea pools are inherited via \texttt{get\_idea\_pool\_for\_ckpt},
and a ranked root-idea selection is recorded by the
\texttt{RootChoice} dataclass so the choice is auditable post hoc
(prompt: \cref{prompt:root_idea_selector}).

Reproducibility hinges on three invariants. (i)~Every randomised
operation seeds itself from the node identifier. (ii)~Every tool
call records its arguments and outputs into the node's working tree,
never the parent's --- this is enforced by the MCP layer's
\texttt{cow\_node\_id} check. (iii)~The cost ledger attaches a
per-node dollar amount at \texttt{CostTracker.init}, so the budget
check is deterministic given the same prompts.

\paragraph{Offline by default.}
Reproducibility also dictates what the search loop is allowed to read.
By default a node's agent loop runs \emph{offline}: the web-access skill
is phase-gated to the paper-writing and replication stages, so a node
cannot consult time-varying live search results during the search itself,
and the per-node trajectory therefore does not depend on the state of the
web on the day it ran. Literature lookup still happens, but earlier and
out of band --- a bounded survey performed at idea time, before the
search begins. A run that genuinely needs live information can opt in
(\code{bfts.allow\_web}) to expose web access to the node agent during
exploration; when it does, the run writes a non-reproducible-trajectory
marker into its checkpoint, so a later reproduction is never silently
treated as exact. The determinism contract stays honest either way: the
default run is byte-reproducible, and a run that trades reproducibility
for live information records that trade in its own provenance.

\subsection{Lineage chaining for computed evidence}\label{sec:exploration:lineage-chaining}

Not every claim a run sets out to support is measured fresh. Some
evidence is \emph{computed} from measurements that already exist --- a
parameter fit over earlier probes, a held-out validation of that fit, a
model-based selection among configurations. Such claims turned out to be
structurally unreachable by independent tree nodes: in a real run,
children directed at fitting or validation but expanded from a parent
that carried no measurements regressed to re-running single probes,
because ``compute from the data you already have'' was never a visible
option --- the child had no data, and nothing told the selector which
node did.

The mechanism that closes this gap uses only the parent$\to$child
working-directory inheritance the tree already has (a child executes
inside a copy of its parent's working tree;
\cref{sec:paper:evidence} relies on the same fact), as one hint per
side of the expansion. On the parent side, the expansion-target
selection prompt (\cref{sec:exploration:node-selection}) is extended
with a coverage hint that names the node whose working tree holds the
most \term{contract-evidence names} --- measurement names required by
the run's declared metric contract, the set of falsifiable claims the
eventual paper must evidence (\cref{sec:paper:contract} defines it in
full) --- so a fit-or-validate child is preferentially expanded from
the node that already holds its inputs. On the child side, a child
whose inherited directory contains lineage measurement files is told,
at start, which files are present and which exact contract names they
contain, and is instructed to compute the derived evidence from those
files rather than re-run the underlying experiments.

Only names cross node boundaries: file names and measurement names,
never values and never a sibling's conclusions. The hints are
run-level coordination metadata of the same kind as the contract
itself, so the branch isolation of \cref{sec:exploration:memory} is
preserved --- a faulty branch's numbers still cannot leak into a
sibling's reasoning --- while the computed-evidence chain
(fit, then validate, then select) can execute along a single lineage
instead of collapsing back into repeated probes.

\subsection{ReAct driver}\label{sec:exploration:react}

\begin{figure}[t]
  \centering
  \resizebox{\columnwidth}{!}{% F04: ReAct loop — Think / Act / Observe / Reflect with numbered stages,
%       paper-style typography.
% style.tikzstyles is loaded by shared/preamble.tex / standalone wrapper.
\begin{tikzpicture}[node distance=2.0cm and 2.5cm]

  \node[state init]                            (think) {\figlabel{react.think}};
  \node[state, right=of think]                 (act)   {\figlabel{react.act}};
  \node[state, below=of act]                   (obs)   {\figlabel{react.observe}};
  \node[state final, below=of think]           (refl)  {\figlabel{react.reflect}};

  % numbered step badges, perched on each state's upper-left corner
  \node[numbered step, above left=0.05cm and -0.15cm of think.north west] {1};
  \node[numbered step, above left=0.05cm and -0.15cm of act.north west]   {2};
  \node[numbered step, above left=0.05cm and -0.15cm of obs.north west]   {3};
  \node[numbered step, above left=0.05cm and -0.15cm of refl.north west]  {4};

  % external annotations (yellow = memory, orange = phase-filtered tools)
  \node[station, fill=ari-orange!22, draw=ari-vermilion,
        above=0.6cm of act]                    (tools) {\figlabel{F4.toolset}};
  \node[station, fill=ari-yellow!30, draw=ari-orange,
        left=0.6cm of think]                   (mem)   {\figlabel{F4.memcontext}};

  % --- big arrows (paper-style) ---
  \draw[big arrow] (think) -- (act)   node[edge label, font=\sffamily\footnotesize]{\figlabel{F4.tool_call}};
  \draw[big arrow] (act)   -- (obs)   node[edge label, font=\sffamily\footnotesize]{\figlabel{F4.execute}};
  \draw[big arrow] (obs)   -- (refl)  node[edge label, font=\sffamily\footnotesize]{\figlabel{F4.parse}};
  \draw[big arrow] (refl)  -- (think) node[edge label, font=\sffamily\footnotesize]{\figlabel{F4.update}};
  \draw[arrow dashed] (obs) to[bend left=18] (think);

  \draw[arrow dashed] (tools.south) -- (act.north);
  \draw[arrow dashed] (mem.east)    -- (think.west);

  % budget annotation
  \node[stage title, below=0.3cm of refl] {\figlabel{F4.maxsteps}};
\end{tikzpicture}
}
  \caption{The ReAct loop. Solid edges are the main cycle; the
           dashed back edge represents an early-exit when the
           observation already closes the open question.}
  \label{fig:react}
\end{figure}

Each node's expansion runs a ReAct-style agent loop in the
\texttt{AgentLoop} class. The maximum step count is
\texttt{MAX\_REACT\_STEPS = 80}; it can be overridden per-instance
via the \texttt{max\_react\_steps} keyword. A separate guard
\texttt{MIN\_TOOL\_CALLS = 2} prevents trivially short trajectories.
The agent's system prompt is reproduced verbatim as
\cref{prompt:agent_system}.

The set of tools available to the agent is filtered per phase by
the tool manager; for example, the exploration phase masks the
paper-writing tools so that BFTS expansion cannot short-circuit
into a draft. A small set of MCP tools is reserved for
\texttt{ari-core} itself (memory CoW operations) and is hidden
from the LLM, ensuring that the model cannot set an arbitrary node
id and bypass the memory skill's copy-on-write check.

\subsection{Research memory}\label{sec:exploration:memory}

Nodes do not run in isolation: each one inherits what its ancestors
established and, in turn, leaves a record for its descendants and for the
paper writer. ARI keeps this record in a per-run \term{research memory}
governed by two requirements a free-form log does not meet. First,
\emph{branch isolation}: a node may read only the memory of its own
ancestors, never that of a sibling or an unrelated checkpoint, so two
parallel leads cannot contaminate each other --- the same ancestor-scope
predicate that defines a lineage in \cref{sec:exploration:lineage}.
Second, \emph{groundability}: because the memory ultimately feeds a
paper, an entry must be able to point at the evidence behind a result
rather than merely restate the result, so a claim can be traced back to
the artefact that supports it.

\subsubsection*{Memory as an index over evidence}

The single source of truth for what a node produced is its structured
self-report --- the \texttt{node\_report.json} written when the node
completes (\cref{sec:paper:evidence} details its fields). A memory entry
does not copy those fields. It carries a short searchable text, a
\term{kind}, a pointer back to the originating node report, and
content-hash references to the specific artefacts --- logs, source files,
output files --- a result rests on. A memory entry is thus an \emph{index
into evidence}, not the evidence itself: the artefacts it cites can be
re-hashed at any time to confirm they still exist and still match, an
integrity check that classifies every reference as verified, missing, or
mismatched. The \term{kind} is drawn from a fixed vocabulary ---
observation, experiment result, failure case, procedure, reflection,
artefact summary, paper claim, and reproducibility event --- so a reader
can ask for exactly the records it needs: the failures along a branch,
say, or only the results that are backed by artefacts.

\subsubsection*{Consolidation at node end}

When a node completes, a deterministic \term{consolidation} step reads
its node report and emits typed entries: an experiment-result entry, with
provenance references to the artefacts it measured, for a node that
produced measurements; a failure-case entry for one that failed. The step
invokes no language model, so the same node report always consolidates to
the same entries --- consolidation is a pure function of recorded state,
and re-running a checkpoint reproduces the memory it produced. It is a
loop hook, not an agent action: the search agent is never asked to ``save
to memory'', and in practice the agent does not reach for the recall
tools on its own. The memory reads the loop depends on are performed by
the loop itself, deterministically --- which is what lets the memory
carry provenance without putting a non-deterministic retrieval on the
critical path (\cref{sec:discussion:p2}).

\begin{figure}[t]
  \centering
  \resizebox{0.92\columnwidth}{!}{% F18: verifiable research memory as a TYPED INDEX OVER EVIDENCE (not a flow).
% The store holds typed records that point, by content hash, at the artefacts
% that ground them; reproduction outcomes are appended and folded to a latest
% status; an admissibility predicate usable(mu) gates which records may become
% paper claims, so the SAME store yields two projections: a working context for
% exploration and a verified context for the paper.
% style.tikzstyles is loaded by shared/preamble.tex / the standalone wrapper.
\begin{tikzpicture}[node distance=1.5cm and 2.2cm]

  % --- source of truth feeding the store -----------------------------
  \node[artifact wide]                          (report) {\figlabel{F18.report}};
  \node[block emph big, below=1.3cm of report]  (store)  {\figlabel{F18.store}};

  % --- what the records index, and what they accrete -----------------
  \node[block native big, left=of store]        (artifacts) {\figlabel{F18.artifacts}};
  \node[block native big, right=of store]       (repro)     {\figlabel{F18.repro}};
  % record vocabulary: a legend in the empty space under the artefacts
  \node[block aux, below=1.3cm of artifacts]    (kinds)     {\figlabel{F18.kinds}};

  % --- two projections of the SAME store -----------------------------
  \node[block native big, below=2.4cm of store]              (workctx) {\figlabel{F18.workctx}};
  \node[guard, below right=1.4cm and 1.7cm of store]         (usable)  {\figlabel{F18.usable}};
  \node[block emph big, below=1.1cm of usable]               (vctx)    {\figlabel{F18.vctx}};

  % --- edges ----------------------------------------------------------
  \draw[arrow]        (report) -- node[edge label] {\figlabel{F18.consolidate}} (store);
  \draw[arrow dashed] (store)  -- node[edge label] {\figlabel{F18.refs}}        (artifacts);
  \draw[arrow]        (repro)  -- node[edge label] {\figlabel{F18.fold}}        (store);
  \draw[arrow dashed] (store)  -- (kinds);
  \draw[arrow]        (store)  -- node[edge label] {\figlabel{F18.forexpl}} (workctx);
  \draw[arrow]        (store)  -- (usable);
  \draw[arrow]        (usable) -- node[edge label] {\figlabel{F18.forpaper}} (vctx);

\end{tikzpicture}
}
  \caption{Verifiable research memory. A deterministic consolidation of
           each node's \texttt{node\_report.json} --- the source of truth
           --- populates a \emph{typed} store that \emph{indexes}
           artefacts (logs, outputs, source) by content hash. Reproduction
           events are appended and folded down to a latest status, and an
           admissibility predicate $\mathrm{usable}(\mu)$ (\cref{eq:usable})
           gates which records reach the verified context. The same store
           yields two readers: the \emph{working context} injected at the
           next node's start (the inheritance path of exploration) and the
           \emph{verified context} that grounds the paper's quantitative
           claims (\cref{sec:paper:evidence}). Ungrounded reflection may
           steer exploration but never the paper body.}
  \label{fig:research-memory}
\end{figure}

\subsubsection*{Two readers: working and verified context}

The store feeds two consumers (\cref{fig:research-memory}), and the
distinction between them is the heart of the design. At the \emph{start}
of every node the loop injects a \term{working context}: the experiment's
fixed core --- goal, target metric, hardware --- together with the
conclusions its ancestors recorded. A child therefore begins from what
its lineage has already established, deterministically and scoped to its
own branch, rather than from an aggregate text dump. At \emph{paper} time
the pipeline builds a \term{verified context} from the best node's
lineage: it folds each result's reproducibility events down to a latest
status and ranks entries so that reproduced results outrank
merely-artefact-grounded ones, which in turn outrank ungrounded
reflection. An entry $\mu$ is \emph{admissible} as a quantitative paper
claim only when it is artefact-grounded and not contradicted by a
reproduction attempt:
\begin{equation}\label{eq:usable}
  \mathrm{usable}(\mu) \;=\;
    \mathrm{grounded}(\mu)\;\wedge\;
    \bigl(\mathrm{repro}(\mu)\neq\mathrm{failed}\bigr).
\end{equation}
Ungrounded reflection can still steer the search through the working
context, but \cref{eq:usable} holds it out of the paper body: the
writer's quantitative claims rest only on admissible, reproduced-first
entries. This is the memory-side counterpart of the evidence assembler
(\cref{sec:paper:evidence}); where the assembler decides which artefact
\emph{bytes} the writer may read, the verified context decides which
\emph{results} are allowed to become claims.

\subsubsection*{Reproducibility as an appended event}

Because past entries are byte-stable --- a node writes only under its own
identifier and existing entries are never edited in place --- a result's
reproducibility status is not patched when the result is later re-run.
Instead the re-run \emph{appends} a reproducibility-event entry, and any
reader folds the events for a target down to their latest state. A result
that fails to reproduce is thereby demoted out of the claim set of
\cref{eq:usable} --- and can be reported honestly as a limitation ---
without rewriting the history that recorded it as promising in the first
place.
